\thispagestyle{empty}

\begin{center}
{\large\bfseries Aplicación Web para la visualización 3D de redes cerebrales \\ \textnormal{\normalsize (Brain Viz)} }\\
\end{center}
\begin{center}
Pablo Moliz Arias\\
\end{center}

%\vspace{0.7cm}

\vspace{0.5cm}
\noindent\textbf{Palabras clave}: \textit{connectoma, red cerebral, visualización 3D, WebGL,
teoría de grafos, matriz de adyacencia, centralidad, React, Three.js, Flask, Docker.}
\vspace{0.7cm}

\noindent\textbf{Resumen}\\
Este Trabajo Fin de Grado aborda el diseño, desarrollo y despliegue de \textbf{Brain Viz}, una
aplicación web para la visualización 3D interactiva de redes cerebrales (\textit{connectomas}).
Estas redes se modelan como grafos, en los que los nodos representan regiones anatómicas del cerebro
y las aristas representan las conexiones entre ellas, y se obtienen a partir de matrices de
adyacencia disponibles en bases de datos públicas de neuroimagen.

Brain Viz sigue una arquitectura cliente-servidor y se estructura en dos componentes desacoplados
que se despliegan juntos en un servidor: un servicio de \textit{backend}, implementado en Python
con Flask y NetworkX, encargado de procesar los ficheros de entrada, construir el grafo, calcular
sus métricas de teoría de grafos y exponer los datos a través de una API; y una aplicación
cliente o \textit{frontend}, desarrollada en React con Three.js mediante \textit{react-three-fiber},
que se ejecuta en el navegador del usuario y es responsable de la visualización 3D interactiva.

La aplicación permite explorar una red filtrando sus conexiones por peso, seleccionar regiones para
consultar sus datos y representar sobre la propia escena las métricas calculadas, tanto las globales
de la red como las centralidades de cada región. Incorpora además cuentas de usuario, con las que un
investigador puede cargar sus propias redes y conservarlas de forma privada, y un rol de
administración que gestiona el catálogo público y los usuarios registrados. El sistema se empaqueta
en contenedores, se apoya en una base de datos para la información persistente y está desplegado en
un servidor público, de manera que se puede utilizar desde cualquier navegador sin instalar nada,
también desde el de un teléfono móvil.

\cleardoublepage

\thispagestyle{empty}

\begin{center}
	{\large\bfseries Web Application for the 3D Visualization of Brain Networks \\ \textnormal{\normalsize (Brain Viz)} }\\
\end{center}
\begin{center}
	Pablo Moliz Arias\\
\end{center}
\vspace{0.5cm}
\noindent\textbf{Keywords}: \textit{connectome, brain network, 3D visualization, WebGL, graph
theory, adjacency matrix, centrality, React, Three.js, Flask, Docker.}
\vspace{0.7cm}

\noindent\textbf{Abstract}\\
This Bachelor's thesis addresses the design, development and deployment of \textbf{Brain Viz}, a web
application for the interactive 3D visualization of brain networks (\textit{connectomes}). These
networks are modelled as graphs, in which nodes represent anatomical brain regions and edges
represent the connections between them, and are obtained from adjacency matrices available in public
neuroimaging databases.

Brain Viz follows a client-server architecture and is structured into two decoupled components that
are deployed together on a server: a \textit{backend} service, implemented in Python with Flask and
NetworkX, responsible for processing the input files, building the graph, computing its graph-theory
metrics and exposing the data through an API; and a client application or \textit{frontend},
developed in React with Three.js via \textit{react-three-fiber}, which runs in the user's browser
and is responsible for the interactive 3D visualization.

The application allows a network to be explored by filtering its connections by weight, selecting
regions to inspect their data and mapping the computed metrics onto the scene itself, both the
global measures of the network and the centrality of each region. It also provides user accounts,
which let a researcher upload their own networks and keep them private, and an administration role
that manages the public catalogue and the registered users. The system is packaged into containers,
relies on a database for persistent information and is deployed on a public server, so that it can
be used from any browser without installing anything, including the browser of a mobile phone.

\cleardoublepage

\thispagestyle{empty}

Yo, Pablo Moliz Arias, alumno de la titulación Ingeniería Informática de la Escuela
Técnica Superior de Ingenierías Informática y de Telecomunicación, con DNI 77448716B,
autorizo la ubicación de la siguiente copia de mi Trabajo Fin de Grado en la biblioteca del
centro para que pueda ser consultada por las personas que lo deseen.

\vspace{1cm}

\noindent Fdo: Pablo Moliz Arias

\vspace{2cm}

\noindent Granada a 25 de junio de 2026.

\cleardoublepage

\thispagestyle{empty}

\noindent\rule[-1ex]{\textwidth}{2pt}\\[4.5ex]

D. \textbf{Juan Ruíz de Miras}, profesor del Departamento de Lenguajes y Sistemas Informáticos.

\vspace{0.5cm}

\textbf{Informa:}

\vspace{0.5cm}

Que el presente trabajo, titulado \textit{\textbf{Aplicación Web para la visualización 3D de redes
cerebrales}}, ha sido realizado bajo su supervisión por \textbf{Pablo Moliz Arias}, y autorizo la
defensa de dicho trabajo ante el tribunal que corresponda.

\vspace{0.5cm}

Y para que conste, expiden y firman el presente informe en Granada a 24 de junio de 2026.

\vspace{1cm}

\textbf{El director: }

\vspace{5cm}

\noindent \textbf{Juan Ruíz de Miras}

% Los agradecimientos se redactarán para la entrega final; se mantienen
% comentados para no dejar una página vacía en las entregas de revisión.
% \chapter*{Agradecimientos}
%
% (Escribe aquí los agradecimientos.)



